% BAN TIENG VIET cua 04_anchor.tex — viet lai 2026-10-12 theo academic-prose.
% Moi con so, cite, label, ref, equation giu NGUYEN so voi ban truoc.
% Thuat ngu theo glossary.tex. KHONG them/bot claim.
\section{Chứng thực neo theo thiết bị}
\label{sec:anchor}

\subsection{Bản ghi của gateway}

Các mệnh đề~\ref{prop:blind}--\ref{prop:coincide} cho thấy tổng lượng và từng độ
sâu vẫn chưa đủ, còn thứ tự thời gian mới đúng là thứ bên gian lận thao túng. Vì vậy, bất kỳ cơ chế kiểm toán nào muốn chặn được tấn công này đều phải quan sát
dấu thời gian, mà bên duy nhất có thể bảo chứng cho một dấu thời gian là thiết bị
đã thực thi lệnh đó. Giả sử gateway giữ một bản ghi
lệnh có chữ ký
\begin{equation}
\mathcal{G}=\{(t_j,u_j,\sigma_j)\},
\label{eq:gateway}
\end{equation}
trong đó:
\begin{symbolist}
\item[$\mathcal{G}$] bản ghi lệnh của gateway, gồm mọi lệnh đóng mở máy bơm trong vụ;
\item[$t_j$] giờ mà bộ điều khiển đóng máy bơm ở lệnh thứ $j$;
\item[$u_j$] độ sâu nước mà bộ điều khiển đã cấp ở lệnh đó, tính bằng mm;
\item[$\sigma_j$] chữ ký số trên cặp $(t_j,u_j)$, do bộ điều khiển tạo ra tại thời điểm đóng mở máy bơm.
\end{symbolist}

Phép thử neo đối chiếu nhật ký khai báo với bản ghi đó bằng \emph{hai hiệu có
hướng} chứ không phải một hiệu đối xứng, vì hai chiều sai lệch khác nhau về bản
chất: một chiều có thể bổ sung từ nguồn, chiều kia thì không.
\begin{equation}
\begin{aligned}
\Delta^{-}&=\{t_j\}\setminus\{t_i\},\\
\Delta^{+}&=\{t_i\}\setminus\{t_j\},\\
\Delta_{\mathrm{depth}}&=\max_{t}\bigl|u(t)-u_{\mathcal{G}}(t)\bigr|,
\end{aligned}
\label{eq:deltas}
\end{equation}
trong đó:
\begin{symbolist}
\item[$\Delta^{-}$] tập giờ có trong bản ghi gateway nhưng vắng trong nhật ký, tức các lệnh bị bỏ sót;
\item[$\Delta^{+}$] tập giờ có trong nhật ký nhưng vắng trong bản ghi gateway, tức các lệnh khai thêm;
\item[$\setminus$] phép hiệu tập hợp, trả về các phần tử thuộc tập thứ nhất mà không thuộc tập thứ hai;
\item[$u(t)$] độ sâu khai báo tại giờ $t$ trong nhật ký $L$, là \emph{tổng} độ sâu của mọi lệnh nằm trong giờ đó theo quy ước đã định nghĩa ở \eqref{eq:uterm}, chứ không phải độ sâu của một lệnh riêng lẻ;
\item[$u_{\mathcal{G}}(t)$] độ sâu mà gateway ghi nhận tại cùng giờ $t$, theo đúng quy ước tổng đó;
\item[$\Delta_{\mathrm{depth}}$] độ lệch lớn nhất giữa hai độ sâu, lấy trên các giờ có mặt ở cả hai bản ghi.
\end{symbolist}
Hai bản ghi đều được lưu theo giờ và mỗi giờ chỉ giữ một mức tổng, nên phép đối
chiếu không làm mất thông tin về số lệnh: mọi lệnh của gateway đã được gộp vào
$u_{\mathcal{G}}(t)$ trước khi so, và chữ ký $\sigma_j$ được kiểm trên \emph{từng}
cặp $(t_j,u_j)$ trước khi bản ghi đó được dùng làm nguồn đối chiếu.

Phán quyết khi đó có ba mức chứ không phải hai:
\begin{equation}
A(L,\mathcal{G})=
\begin{cases}
\text{giả mạo}, & \Delta^{+}\neq\emptyset\ \vee\ \Delta_{\mathrm{depth}}>\epsilon,\\
\text{đối soát}, & \Delta^{+}=\emptyset\ \wedge\ \Delta^{-}\neq\emptyset,\\
\text{khớp}, & \text{trường hợp còn lại},
\end{cases}
\label{eq:anchor}
\end{equation}
trong đó:
\begin{symbolist}
\item[$A(L,\mathcal{G})$] phán quyết của phép thử neo đối với nhật ký $L$;
\item[$\epsilon$] dung sai so khớp độ sâu, bằng $10^{-6}$~mm, tức khớp chính xác tới mức nhiễu dấu phẩy động;
\item[$\vee$] phép hoặc logic, nghĩa là chỉ cần khai thêm một lệnh hoặc lệch một độ sâu là đủ kết luận giả mạo;
\item[$\wedge$] phép và logic;
\item[\text{đối soát}] cờ đánh dấu nhật ký thiếu lệnh so với gateway; đây là \emph{lỗi ghi chép} chứ không phải kết luận gian lận, và các mục thiếu được bổ sung từ $\mathcal{G}$.
\end{symbolist}
Gọi tên riêng ba mức này là cần thiết, vì mỗi lớp gian lận để lại dấu vết ở một
chỗ khác nhau: dịch thời điểm và thêm lệnh đều tạo ra giờ không có trong
$\mathcal{G}$ nên làm $\Delta^{+}$ khác rỗng, thổi phồng độ sâu chỉ làm thay đổi
$\Delta_{\mathrm{depth}}$, còn bỏ sót lệnh chỉ làm $\Delta^{-}$ khác rỗng và do đó
chỉ kích hoạt cờ đối soát. Nhờ vậy một phép đối chiếu duy nhất vẫn bắt được mọi
lớp gian lận, nhưng không đánh đồng việc quên ghi với việc khai thêm.

\begin{proposition}[Tính đầy đủ của neo thiết bị]
\label{prop:complete}
Nếu nhật ký khai báo $L$ khai một giờ không có trong bản ghi gateway
$\mathcal{G}$, hoặc khai sai độ sâu tại một giờ chung, thì $A(L,\mathcal{G})$ trả
về \text{giả mạo}. Nếu $L$ chỉ thiếu những lệnh mà $\mathcal{G}$ có, thì
$A(L,\mathcal{G})$ trả về \text{đối soát}: các mục thiếu được bổ sung từ
$\mathcal{G}$ và không dùng làm căn cứ kết tội. Nếu $L$ và $\mathcal{G}$ khớp cả
tập giờ lẫn độ sâu, thì $L$ là một bản sao trung thực và không có sự giả mạo nào
đối với ba đại lượng trong \eqref{eq:three}.
\end{proposition}

\begin{proof}
Xem Phụ lục~S12. Hai mệnh đề đầu là định nghĩa của \eqref{eq:anchor}, còn mệnh đề
cuối suy ra được vì cả ba đại lượng có thể làm giả đều là hàm của tập giờ và các
độ sâu theo giờ.
\end{proof}

\subsection{Vì sao neo thiết bị không dẫn tới kết tội oan}

Điểm hoá giải là neo thiết bị đối chiếu một bản sao với \emph{nguồn} chứ không đối chiếu
một nhật ký với một \emph{mô hình}: nguồn sự thật là bản ghi của gateway chứ không phải
trí nhớ của nông dân, nên nhật ký chỉ là một góc nhìn của bản ghi mà thiết bị đã giữ. Vì
vậy \eqref{eq:anchor} tách hai chiều sai lệch: nếu nhật ký bỏ sót một lệnh có trong
$\mathcal{G}$ thì đó là lỗi ghi chép, mục đó được bổ sung từ $\mathcal{G}$ và phán quyết
chỉ dừng ở mức \text{đối soát}; còn nếu nhật ký khai thêm một lệnh không có trong
$\mathcal{G}$, hoặc khai sai thời điểm hay độ sâu so với $\mathcal{G}$, thì đó là
\text{giả mạo}. Khi đó phép thử neo chỉ có thể kết tội oan một nông dân trung thực nếu bản ghi
của gateway tự mâu thuẫn với chính nó, và đó là bài toán bảo đảm toàn vẹn thiết bị chứ
không còn là bài toán kiểm toán nhật ký. Nhờ vậy Bảng~\ref{tab:falseacc} giữ ở mức
$0/60$, trong khi các bộ phát hiện dựa trên quỹ đạo phát lại vẫn kết tội oan ở mức
$26{,}7\%$ và $5{,}0\%$ (Phụ lục~S18).

\subsection{Chi phí và phạm vi của tuyên bố}

Neo thiết bị yêu cầu dùng gateway trong bản thảo chưa công bố nêu trên~\cite{twingate2026}, thiết bị vốn đóng mở máy bơm và ghi lại bản ghi lệnh có chữ ký kèm dấu thời gian
như một phần của chức năng an toàn. Chi phí biên của việc chứng thực chỉ gồm lưu trữ và truyền đi một bản ghi vốn đã được gateway sinh ra trong quá trình hoạt động bình thường. Cũng cần phân biệt rõ phạm vi của bằng chứng này: một \emph{lệnh bơm có chữ ký} chứng thực rằng bộ điều khiển đã ban hành lệnh tại một thời điểm xác định, còn nó không chứng thực \emph{lượng nước thực sự đã được cấp}, vì bài này không dùng phép đo lượng nước độc lập nào ngoài bản ghi đó (Bảng~S4 của Phụ lục).

Tuy vậy, phạm vi áp dụng của tuyên bố hẹp hơn kỳ vọng mà phần mở đầu gợi ra. Một nhật ký viết tay, không có thiết bị nào tham gia vào vòng ghi nhận, thì bản
thân nó \emph{không} tạo thành bằng chứng được chấp nhận, vì
Mục~\ref{sec:results-attack} đo được một bên gian lận duy lý có thể thổi phồng tín chỉ
thêm $+35\%$ mà vẫn lọt qua mọi phép thử dựa trên vật lý và tổng lượng ($6/6$
nhật ký bị sửa). Vì vậy MRV với chi phí \emph{biên} gần như bằng không chỉ
khả thi ở nơi có một bộ điều khiển đứng giữa nông dân và máy bơm, còn ở nơi tưới
thủ công thì lớp này không áp dụng được.
